% Irodov Part III §3.5 last 275--287; PDF pp.141--143,
% printed pp.132--134. Graph of original Fig 3.76 consulted.
\electromagproblem{275}{132}
一根很长的圆柱形顺磁导线
流轴向总电流 $I$，
均匀磁化率为 $\chi$。
求：(a) 导线表面束缚分子电流总量；
(b) 体内束缚分子电流总量；
并说明相对流向。
\exsolution
导线半径记为 $R$。
自由电流密度
$j=I/(\pi R^2)$，
线内
$H_\varphi(r)=jr/2=Ir/(2\pi R^2)$。
磁化强度
$J_\varphi(r)=\chi H_\varphi(r)
=\chi I r/(2\pi R^2)$。
体束缚电流密度
$\bm j_b=\nabla\times\bm J$，
其轴向分量
\[
j_{b,z}=\frac1r\frac{\mathrm d}{\mathrm dr}
(rJ_\varphi)
=\frac{\chi I}{\pi R^2}.
\]
截面积分体束缚电流
\[
\boxed{I'_{\rm vol}=\chi I}.
\]
表面束缚电流密度
$\bm K_b=\bm J\times\bm n$，
在 $r=R$，
$\bm J\parallel\bm e_\varphi$、
$\bm n=\bm e_r$，
故 $\bm K_b\parallel-\bm e_z$，
沿周长积分模为
\[
\boxed{I'_{\rm surf}=
2\pi R\,J_\varphi(R)=\chi I}.
\]
\excheck 两者\emph{大小相等、方向相反}，
导线及外侧整体的束缚电流净量为零。

\electromagproblem{276}{132}
图3.75：很长的通电螺线管
从中部 $O$ 开始，轴线 $x>0$
的半段填充相对磁导率为 $\mu$ 的介质，
$x<0$ 为真空。
画出轴线上 $B(x)$、$H(x)$、
磁化强度 $J(x)$ 的大致分布。
\exsolution
远离分界面，螺线管的自由表面电流
等效密度为 $K=nI$。
真空一侧（$x\ll-R$）
\[
\boxed{H_-\simeq nI,\qquad
B_-\simeq\mu_0nI,\qquad J_-=0}.
\]
介质一侧（$x\gg R$）
\[
\boxed{H_+\simeq nI,\qquad
B_+\simeq\mu\mu_0nI,\qquad
J_+\simeq(\mu-1)nI}.
\]
在 $O$ 附近，由端部和磁荷退极化效应，
$B$、$H$ 的真实变化不能只以
两边常值拼接；沿轴无自由磁荷，
磁感应强度的法向分量连续，
故 $B(x)$ 应平滑连接两侧平台。
磁化强度在介质外为零，
进入介质后建立并趋于平台。
\excheck 此处说明的是原书要求的
\emph{定性图形}，不是磁介质有限端部精确边值解。
若只把 $\mu$ 作为乘数应用到整个轴上，
会错误地改变真空侧远处的场。

\electromagproblem{277}{132}
无限长直导线通电流 $I$，
恰位于两种非导电磁介质的平面交界线上，
相对磁导率分别为 $\mu_1,\mu_2$。
磁力线为以导线为中心的圆周，
求距导线 $r$ 处磁感应强度模。
\exsolution
在同一环路上，$\bm B$
沿界面切向的法向跨界条件以及
题设圆形磁力线使其模在两半圆上相同。
相应磁场强度分别为
$H_1=B/(\mu_0\mu_1)$、
$H_2=B/(\mu_0\mu_2)$。
半径为 $r$ 的整圆安培环路由两半圆组成，
\[
I=\oint\bm H\cdot\mathrm d\bm\ell
=\pi r\frac B{\mu_0\mu_1}
+\pi r\frac B{\mu_0\mu_2}.
\]
解得
\[
\boxed{B(r)=\frac{\mu_0I\,\mu_1\mu_2}
{\pi r(\mu_1+\mu_2)}}.
\]
\excheck $\mu_1=\mu_2=1$ 时
恢复真空直导线场
$B=\mu_0I/(2\pi r)$。

\electromagproblem{278}{132}
半径为 $R$ 的环形电流位于
真空与磁介质的\emph{平面边界内}，
磁介质相对磁导率 $\mu$。
若无磁介质时场为 $\bm B_0$，
求该构型在电流环轴上及其余空间的磁感应强度。
\exsolution
把电流环所在平面作为介质分界面，
两侧均是无自由体电流区域。
对各个毕奥—萨伐尔自由电流元产生的
磁矢势使用线性材料边界条件，
可将本几何中的总场写成
原均匀真空问题的固定比例：
\[
\bm B(\bm r)=\frac{2\mu}{1+\mu}\bm B_0(\bm r).
\]
因此轴上及其它适用的场点：
\[
\boxed{B=\frac{2\mu}{1+\mu}B_0}.
\]
\excheck $\mu=1$ 还原 $\bm B_0$；
$\mu\to\infty$ 比例趋于2。
该结果依赖环电流\emph{严格位于两介质平界面}
的特殊几何，而非任意介质分布都适用。

\electromagproblem{279}{132}
相对磁导率为 $\mu$ 的
均匀各向同性磁介质球
置于外加均匀磁场 $\bm B_0$ 中，
球内均匀磁化。
已知均匀磁化球的退磁场
$\bm H_d=-\bm J/3$。
求球内磁感应强度。
\exsolution
外加磁场强度
$\bm H_0=\bm B_0/\mu_0$，
球内
$\bm H=\bm H_0-\bm J/3$。
线性材料磁化
$\bm J=(\mu-1)\bm H$，
因此
\[
\bm H=\bm H_0-\frac{\mu-1}{3}\bm H,
\quad
\bm H=\frac3{\mu+2}\bm H_0.
\]
材料内
$\bm B=\mu\mu_0\bm H$，所以
\[
\boxed{\bm B=
\frac{3\mu}{\mu+2}\bm B_0}.
\]
\excheck $\mu\to1$ 时系数为1；
$\mu\to\infty$ 时系数趋于3。

\electromagproblem{280}{132}
长 $\ell=15\,\mathrm{cm}$ 的
圆柱形永久磁铁上均匀绕
$N=300$ 匝线圈。
通电流 $I=3.0\,\mathrm A$ 后
磁铁外磁场恰好消失，
求磁性材料的矫顽力 $H_c$。
\exsolution
线圈在长圆柱轴上产生磁场强度
$H_{\rm coil}\simeq nI=NI/\ell$。
要抵消永久磁化所产生的外场，
需施加相反方向的矫顽磁场。
故矫顽力模为
\[
\boxed{H_c=\frac{NI}{\ell}}.
\]
代数值：
\[
\boxed{H_c=\frac{(300)(3.0)}{0.15}
=6.0\times10^3\,\mathrm{A/m}
=6.0\,\mathrm{kA/m}}.
\]
\excheck 题目测量使用的是磁场强度
$\bm H$ 而不是磁感应强度 $\bm B$，
单位为 $\mathrm{A/m}$。

\electromagproblem{281}{132}
永久磁铁制成平均直径
$d=20\,\mathrm{cm}$ 的闭合圆环，
环上留宽
$b=2.0\,\mathrm{mm}$ 的小气隙。
气隙内磁感应强度
$B=40\,\mathrm{mT}$。
忽略边缘漏磁，求磁铁内部磁场强度模。
\exsolution
没有自由绕组电流，
沿闭合环路有
$\oint\bm H\cdot\mathrm d\bm\ell=0$。
设铁磁体中的磁场强度为 $H_m$，
气隙中
$H_g=B/\mu_0$。
忽略气隙长度对主圆周的微小影响：
\[
H_m(\pi d-b)+\frac B{\mu_0}b=0.
\]
因此
\[
\boxed{|H_m|=
\frac{bB}{\mu_0(\pi d-b)}
\simeq\frac{bB}{\mu_0\pi d}}.
\]
代入：
\[
\boxed{|H_m|\simeq1.0\times10^2\,
\mathrm{A/m}=0.10\,\mathrm{kA/m}}.
\]
\excheck 磁铁内部的 $\bm H_m$
与气隙场方向相反；
此处报磁场强度的模。

\electromagproblem{282}{132--133}
平均半径
$R=250\,\mathrm{mm}$ 的环形铁芯，
绕有 $N=1000$ 匝，
电流 $I=0.85\,\mathrm A$。
铁芯切有气隙
$b=1.00\,\mathrm{mm}$，
气隙磁感应强度
$B=0.75\,\mathrm T$。
忽略漏磁，求该工况下铁的相对磁导率 $\mu$。
\exsolution
沿环路应用安培定律：
\[
NI=H_{\rm iron}(2\pi R-b)+H_{\rm gap}b.
\]
两处磁通守恒且截面相同，
$B_{\rm iron}\simeq B_{\rm gap}=B$。
因此
$H_{\rm iron}=B/(\mu\mu_0)$、
$H_{\rm gap}=B/\mu_0$。
解出
\[
\boxed{\mu=
\frac{B(2\pi R-b)}
{\mu_0NI-Bb}}.
\]
在 $b\ll R$ 下近似
$\mu\simeq2\pi RB/(\mu_0NI-Bb)$。
代入：
\[
\boxed{\mu\simeq3.7\times10^3}.
\]
\excheck 分母 $\mu_0NI-Bb$ 必须为正，
否则所假设的同向磁化与场分布
不适用于题设工况。

\electromagproblem{283}{133}
原书图3.76给出了商业纯铁的初始磁化曲线
$B(H)$。由图推绘相对磁导率 $\mu(H)$，
并估计最大值与对应磁场强度。
\exsolution
在各个图上的非零 $H$ 点，
相对磁导率按定义为
\[
\boxed{\mu(H)=\frac{B(H)}{\mu_0H}}.
\]
因此应逐点用纵坐标 $B$ 除以
横坐标 $H$（以 SI
$\mathrm{A/m}$ 代入），再把所得点作图。
在 $H$ 很小时，$B$ 起初随 $H$ 迅速增加，
$\mu$ 上升；到饱和区
$B$ 增长缓慢，而分母 $H$ 增大，
故 $\mu$ 下降。
从原书图3.76读数，
最大值大约位于
\[
\boxed{H_{\max}\simeq0.06\,\mathrm{kA/m}
=60\,\mathrm{A/m}},
\]
该处磁感应强度约 $0.75\,\mathrm T$，
所以
\[
\boxed{\mu_{\max}\approx
\frac{0.75}{(4\pi\times10^{-7})(60)}
\simeq1.0\times10^4}.
\]
\excheck 此题采用\emph{从印刷图读数}
的近似值，不具有精密物性测量的有效数字；
$\mu(H)$ 是非线性材料的割线磁导率。

\electromagproblem{284}{133}
图3.76同一铁磁材料，
平均直径
$d=50\,\mathrm{cm}$ 的细环形铁芯
绕 $N=800$ 匝，
电流 $I=3.0\,\mathrm A$，
铁芯气隙宽
$b=2.0\,\mathrm{mm}$。
忽略漏磁，结合原图求工作点 $\mu$。
\exsolution
沿磁路
$NI=H(\pi d-b)+Hb_g$，
其中气隙内
$H_g=B/\mu_0$，
因此
\[
B=\frac{\mu_0NI}{b}
-\frac{\mu_0(\pi d-b)}{b}\,H.
\]
代入数据并把 $H$ 的横坐标写作
$\mathrm{kA/m}$：
\[
\boxed{B(\mathrm T)\simeq
1.51-0.987H(\mathrm{kA/m})}.
\]
与原书图3.76的非线性磁化曲线求交，
读得
$H\simeq0.26\,\mathrm{kA/m}$、
$B\simeq1.25\,\mathrm T$。
相对磁导率
\[
\boxed{\mu=\frac B{\mu_0H}
\simeq\frac{1.25}
{(4\pi\times10^{-7})(260)}
\simeq3.8\times10^3\approx4\times10^3}.
\]
\excheck \textbf{书后答案数量级冲突：}
书后简答似给“$4\times10^2$”，
但使用同一书后读数 $B\approx1.25$T、
$H\approx0.26$kA/m，
必得 $\sim4\times10^3$，
相差10倍；本稿按 SI 数值计算。

\electromagproblem{285}{133}
横截面积 $S$、磁化率为 $\chi$
的细长顺磁棒，
沿通电线圈轴向放置。
一端位于线圈中心、局部磁感应强度为 $B$，
另一端位于磁场可以忽略的远处。
求线圈对棒的磁力。
\exsolution
弱顺磁体的磁化
$J=\chi H\simeq\chi B/\mu_0$。
恒定线圈电流下，
棒单位体积的可用磁势能近似
$u=-\chi B^2/(2\mu_0)$。
棒沿轴微量进入线圈，
相当于有体积 $S\,\mathrm dx$
从无场处移到场强 $B$ 处。
有效能下降
$\mathrm dU\simeq
-[\chi B^2/(2\mu_0)]S\,\mathrm dx$，
吸入力模为
\[
\boxed{F\simeq\frac{\chi SB^2}{2\mu_0}}.
\]
\excheck 方向为将顺磁棒\emph{吸入}
强磁场一侧；
该式只适用于弱磁化与细棒近似。

\electromagproblem{286}{133--134}
图3.77中体积
$V=41\,\mathrm{mm^3}$
的顺磁小球位于电磁铁极靴轴线上。
磁场随高度变化为
$B(x)=B_0e^{-a x^2}$，
$B_0=1.50\,\mathrm T$、
$a=100\,\mathrm{m^{-2}}$。
求 (a) 磁吸引力最大的高度；
(b) 若 $F_{\max}=160\,\mathrm{\mu N}$，
求材料磁化率 $\chi$。
\exsolution
弱顺磁小球的磁势能近似为
$U(x)=-\chi VB(x)^2/(2\mu_0)$。
磁力沿 $x$ 方向
$F_x=-\mathrm dU/\mathrm dx
=[\chi V/(2\mu_0)]\,\mathrm d(B^2)/\mathrm dx$。
因为
$B^2=B_0^2e^{-2ax^2}$，
吸引力模
\[
|F(x)|=\frac{2\chi Va B_0^2}{\mu_0}
x\,e^{-2ax^2},\quad x>0.
\]
对 $x e^{-2ax^2}$ 求导，
极值条件
$1-4ax^2=0$，于是
\[
\boxed{x_{\max}=\frac1{2\sqrt a}
=0.050\,\mathrm m=5.0\,\mathrm{cm}}.
\]
在该处
$|F_{\max}|=
\chi VB_0^2\sqrt a\,e^{-1/2}/\mu_0$，
所以
\[
\boxed{\chi=
\frac{\mu_0F_{\max}e^{1/2}}
{VB_0^2\sqrt a}
\simeq3.6\times10^{-4}}.
\]
\excheck 球体积换算
$41\,\mathrm{mm^3}=41\times10^{-9}\,\mathrm{m^3}$。
书后纯文本提取的 $x_m$ 可能丢平方根；
从极值方程所得长度尺度必须是 $1/\sqrt a$。

\electromagproblem{287}{134}
一颗体积为 $V$、
顺磁磁化率为 $\chi$ 的小球，
从线圈轴上磁感应强度为 $B$ 的位置，
被缓慢移动到场可以忽略的远处。
求外力所做功。
\exsolution
弱线性顺磁球在给定外场中的
等效磁势能（保持线圈电流不变）
为
$U=-\chi VB^2/(2\mu_0)$。
初态
$U_i=-\chi VB^2/(2\mu_0)$，
末态
$U_f\simeq0$。
准静态移走小球的外力功
$A_{\rm ext}=U_f-U_i$：
\[
\boxed{A_{\rm ext}\simeq
\frac{\chi VB^2}{2\mu_0}}.
\]
\excheck 对顺磁材料 $\chi>0$，
外界需克服吸引力做正功；
若改为抗磁材料 $\chi<0$，
力方向与功号将相反。
